% =============================================================================
% frontmatter.tex  —  Title page, Abstract, How to Read This Manuscript
% Part of the K-Bound long-form monograph.
% NO \documentclass or \usepackage here; this file is \input by main.tex.
% =============================================================================

% ---------------------------------------------------------------------------
% Title page
% ---------------------------------------------------------------------------
\begin{titlepage}
\centering
\vspace*{2.5cm}

{\Huge\bfseries When Is Label-Free Adaptation Knowable?}\\[1.2em]
{\large\itshape The Knowability Boundary of Label-Free Adaptation --- Manuscript Edition}

\vspace{2.5cm}

{\LARGE Pratik Niroula}\\[0.6em]
{\large Independent Researcher}

\vspace{3cm}

{\large 2026}

\vspace{3.5cm}

\begin{minipage}{0.78\textwidth}
\small\itshape
Long-form companion to the working paper ``When Is Label-Free
Adaptation Knowable?''\ (K-Bound, v0.4).
Every theorem, number, table, and figure in this manuscript reproduces
from the public repository at
\url{https://github.com/Pratikn03/AutoML_Flagship_V8}.
Nothing has been added beyond the exposition; nothing has been rounded,
selected, or otherwise adjusted from the verified artifact manifests.
\end{minipage}

\vfill
\end{titlepage}

% ---------------------------------------------------------------------------
% Abstract
% ---------------------------------------------------------------------------
\chapter*{Abstract}
\addcontentsline{toc}{chapter}{Abstract}

Test-time adaptation (TTA) improves models on unlabeled target data, yet the
very unlabeled objective that recovers accuracy under one distribution shift can
silently degrade it under another --- and, without labels, a deployed system
cannot tell which is happening.
Existing methods propose increasingly robust \emph{adaptation mechanisms}; we
instead study a prior question: \emph{can a system decide, without labels,
whether adaptation should happen at all?}

We formalize adaptation benefit as the excess risk between an adapted candidate
and a frozen baseline conditioned on label-free observable evidence, and separate
test-time regimes into \emph{knowably helpful}, \emph{knowably harmful}, and
\emph{unknowable}.
We prove (Theorem~\ref{thm:imp}, with an explicit witness) that when two target
worlds induce the same observable evidence but opposite benefit, no label-free rule
can certify the correct decision in both, so the worst-case-optimal action on such
evidence is to abstain.
Under an observable risk-alignment assumption we give a finite-sample
adapt/freeze/abstain certificate (Theorem~\ref{thm:cert}) that controls the
false-adapt and false-freeze rates at a chosen level, and we identify a provable
positive regime (Theorem~\ref{thm:pos}, covariate shift).
A quantitative strengthening via the Le~Cam two-point lower bound
(Theorem~\ref{thm:imp}) makes the impossibility result a tight minimax statement:
the optimal committal error equals $1 - \TV$ of the evidence distributions, which
is~$1$ in the exact witness and decays to zero only as the evidence separates.

We instantiate the framework (\KGA) on a 123-task anomaly-routing benchmark and a
controlled mixed-regime study.
\KGA\ abstains where the true benefit is near zero, refuses a fusion path that
hurts on~80\% of tasks (cutting regret-to-oracle $\approx 11{\times}$ versus
always-adapt), and in the mixed regime beats an always-freeze policy by
$\approx 0.10$ AUROC.
In \emph{online non-stationary} CIFAR-10 test-time adaptation, across a
helpful-dominated and a harm-dominated stream, no fixed policy is robust:
always-adapt collapses in the harm regime ($0.339$) and always-freeze forgoes gains
in the helpful regime --- while \KGA\ is near-best in both, giving the highest mean
accuracy across regimes ($0.490$ vs.\ $0.456$ / $0.444$).
On a \emph{pre-registered} CIFAR-10-C stress grid (severity $\times$ batch $\times$
composition $\times$ aggressiveness; harmful base rate $0.16$--$0.34$), \KGA\
\emph{beats both} trivial policies for Tent and EATA --- regret-to-oracle $0.0016$
versus always-adapt $0.009$ and always-freeze $0.12$, at a $0\%$ false-adapt
rate --- and ties the collapse-resistant SAR.

We are explicit about scope: where harmful adaptation is mild, always-adapt is
already strong, so the certificate's distinctive value is bounded by how often
\emph{catastrophic, detectable} harm occurs --- itself the quantity the theory
characterizes.
The surviving open piece is the label-free \emph{bracketing} of the ordinal
accuracy comparison on the observable disagreement region in the multiclass and
regression cases (Conjecture~\ref{conj:gen}); every other element of the
characterization is proved.

% ---------------------------------------------------------------------------
% How to Read This Manuscript
% ---------------------------------------------------------------------------
\chapter*{How to Read This Manuscript}
\addcontentsline{toc}{chapter}{How to Read This Manuscript}

This monograph presents the K-Bound research program at a depth that the
20-page working paper cannot sustain.
Theorems and experiments are stated faithfully from the paper; the additional
length is entirely expository --- motivation, worked examples, proof commentary,
and explicit connections between the abstract theory and its empirical realizations.
A reader who has the paper in hand can use the chapter headings as an index;
a reader coming to K-Bound fresh should read Chapters~\ref{ch:intro}
and~\ref{ch:related} first to orient on the problem and its neighbours, then
follow Chapters~\ref{ch:formulation} through~\ref{ch:theory2} for the theory,
and Chapters~\ref{ch:kga} through~\ref{ch:exp2} for the instantiation and
experiments.

\medskip
\noindent\textbf{Chapter~\ref{ch:intro}: Introduction.}
Sets up the TTA double-edge problem and the central question of knowability.
Introduces the adapt/freeze/abstain trichotomy, previews all five theoretical
results and the headline experimental numbers, and gives an honest statement of
scope.

\medskip
\noindent\textbf{Chapter~\ref{ch:related}: Background and Related Work.}
A thorough review of the five lines of work adjacent to K-Bound: TTA mechanisms,
TTA failure and collapse studies, guarded/protected adaptation, label-free
accuracy estimation, unsupervised risk and identifiability, selective prediction
and abstention, conformal and anytime-valid inference, and covariate-shift theory.
Closes with a positioning table and a precise statement of the gap K-Bound fills.

\medskip
\noindent\textbf{Chapter~\ref{ch:formulation}: Problem Formulation.}
The formal setup: source and target distributions, the adaptation benefit
$\Delta = R_T(f_0) - R_T(f_a)$, observable evidence $Z$, the three regimes, and
the key definitions of risk alignment and knowability.
Includes the ``label-free'' clarification box and the knowability phase diagram.

\medskip
\noindent\textbf{Chapter~\ref{ch:theory1}: Impossibility and the Unknowable Regime.}
The first two theoretical results.
Theorem~\ref{thm:imp} (non-identifiability with explicit witness) and its Le~Cam
quantitative strengthening establish the unknowable regime as fundamental.
Theorem~\ref{thm:gate} gives the Bayes-optimal gate and the plug-in regret
decomposition, which justifies abstaining in the low-margin band.

\medskip
\noindent\textbf{Chapter~\ref{ch:theory2}: Certificates and Positive Regimes.}
The remaining three theoretical results.
Theorem~\ref{thm:cert} is the finite-sample adapt/freeze/abstain certificate with
a controlled false-adapt rate; Corollaries give forced abstention and sample
complexity.
Theorem~\ref{thm:pos} proves a knowable positive regime under covariate shift.
Theorems~\ref{thm:disagree}, \ref{thm:disagree-mc}, and~\ref{thm:disagree-reg}
reduce sign-of-difference identifiability to an ordinal accuracy comparison on the
observable disagreement region for binary, multiclass, and regression losses.
Conjecture~\ref{conj:gen} states the remaining open piece.

\medskip
\noindent\textbf{Chapter~\ref{ch:kga}: The \KGA\ Algorithm.}
The \KGA\ (Knowability-Guided Adaptation) wrapper: evidence extraction,
benefit estimation, conformal radius, and the three-way decision.
Assumptions, computational cost, and the relationship to any underlying TTA
candidate.

\medskip
\noindent\textbf{Chapter~\ref{ch:exp1}: Experiments I --- Anomaly Routing.}
The 123-task clean suite, the harmful-regime refusal experiment, the controlled
mixed-regime study, and the non-identifiability witness (both empirical and exact).
Evidence ablations and the $\alpha$-sweep safety/coverage tradeoff.

\medskip
\noindent\textbf{Chapter~\ref{ch:exp2}: Experiments II --- Online TTA and CIFAR-10-C.}
The online non-stationary CIFAR-10 cross-regime test, the per-corruption CIFAR-10-C
helpful-dominated protocol, and the pre-registered CIFAR-10-C stress grid.
Breadth results on 62 additional harmful-dominated label-free tasks.

\medskip
\noindent\textbf{Chapter~\ref{ch:elara}: The ELARA/RGA Instantiation.}
Reliability-gated multimodal fusion as the worked, code-validated special case:
the ten-result ELARA theorem stack (T1--T9, GDR) and its mapping to the general
K-Bound theorems.

\medskip
\noindent\textbf{Chapter~\ref{ch:discussion}: Discussion, Limitations, and Conclusion.}
When to deploy K-Bound; honest scope; limitations; and the research frontier
(ImageNet-C / ViT scale, natural-shift benchmarks, the open conjecture).

\medskip
\noindent\textbf{Appendix~\ref{app:repro}: Reproducibility.}
One-command reproduction; data provenance and checksums; environment and
determinism notes.

\medskip
\noindent\textbf{Appendix~\ref{app:tables}: Extended Tables.}
The full 65-cell CIFAR-10-C per-corruption-by-severity breakdown and
auxiliary result tables referenced from the main chapters.

\medskip
\noindent\textbf{Appendix~\ref{app:notation}: Notation Reference.}
A single-page glossary of all symbols used in the monograph.

\medskip\noindent
Readers interested primarily in the \emph{theory} can read
Chapters~\ref{ch:formulation}--\ref{ch:theory2} as a self-contained unit.
Readers interested primarily in the \emph{practical algorithm} can read
Chapter~\ref{ch:kga} with only the informal summaries of
Theorems~\ref{thm:imp} and~\ref{thm:cert}.
Both paths converge at Chapter~\ref{ch:discussion}, which synthesizes the
implications for deployment.
